% ============================================================================
% PATENTE DE INVENÇÃO — Ψ-NET PROTOCOL
% Instituto Nacional da Propriedade Industrial (INPI) — Brasil
% Natureza: Patente de Invenção (PI)
% Classificação IPC: H04L 69/22; H04L 9/00; G06Q 20/06; H04L 45/00
% ============================================================================

\documentclass[12pt,a4paper]{article}

\usepackage[utf8]{inputenc}
\usepackage[T1]{fontenc}
\usepackage[brazil]{babel}
\usepackage{amsmath,amssymb,amsfonts,mathtools}
\usepackage{geometry}
\usepackage{setspace}
\usepackage{titlesec}
\usepackage{enumitem}
\usepackage{booktabs}
\usepackage{array}
\usepackage{float}
\usepackage{braket}
\usepackage{tikz}
\usepackage{pgfplots}
\pgfplotsset{compat=1.18}
\usetikzlibrary{shapes.geometric,arrows.meta,positioning,calc,
  decorations.pathreplacing,patterns,fit,backgrounds,
  decorations.markings,3d,chains}

\geometry{left=3cm, right=2cm, top=3cm, bottom=2cm}
\onehalfspacing

\titleformat{\section}{\normalfont\Large\bfseries}{\thesection}{1em}{}
\titleformat{\subsection}{\normalfont\large\bfseries}{\thesubsection}{1em}{}
\titleformat{\subsubsection}{\normalfont\normalsize\bfseries}{\thesubsubsection}{1em}{}

% Comandos auxiliares
\newcommand{\alphaii}{\alpha^2}
\newcommand{\Dfolds}{D_{\text{folds}}}
\newcommand{\Fexp}{F_{\text{exp}}}
\newcommand{\thetaM}{\theta_{\text{Miguel}}}
\newcommand{\Hfour}{\mathcal{H}_4}
\newcommand{\PsiBit}{\textsc{PsiBit}}
\newcommand{\PsiNet}{\Psi\text{-NET}}

\begin{document}

\begin{center}
{\Large\bfseries RELATÓRIO DESCRITIVO}\\[0.8cm]
{\large\bfseries MÉTODO E SISTEMA DE PROTOCOLO DE COMUNICAÇÃO\\
EM REDE BASEADO EM PILHA LUMINODINÂMICA DE\\
CINCO CAMADAS COM ENDEREÇAMENTO ANGULAR,\\
ROTEAMENTO HOLOGRÁFICO E CONSENSO POR PROVA DE FASE}\\[0.8cm]
\textbf{Inventor:} Luiz Antonio Rotoli Miguel\\[0.3cm]
\textbf{Classificação IPC:} H04L 69/22; H04L 9/00; G06Q 20/06; H04L 45/00\\[0.3cm]
\textbf{Referência a pedidos anteriores:}\\
BR 10 2025 026951 1 (ACOM --- Compressão Ontológica de Memória)\\
ACOM Mirror (Reflexão Dimensional de Dados)\\
ACOM Mirror Crypto (Criptografia por Separação Ontológica)
\end{center}

\vspace{1cm}

% ============================================================================
\section{Campo da invenção}
% ============================================================================

A presente invenção refere-se a métodos e sistemas de protocolo
de comunicação em rede baseados em pilha luminodinâmica de cinco
camadas, aplicável aos setores de telecomunicações, redes de
computadores, criptomoedas, segurança da informação e Internet.
Mais especificamente, a invenção descreve um protocolo de
comunicação denominado $\PsiNet$ que substitui a pilha de
protocolos convencional (modelo OSI/TCP-IP) por uma arquitetura
de cinco camadas fundamentada na Teoria da Gravitação
Luminodinâmica (TGL), compreendendo: Camada~0 (Enlace Psiônico)
com unidade fundamental \PsiBit{} de 2 bits isomórfica à
ligação psiônica; Camada~1 (Transporte Radical) baseada na
radicalização $g = \sqrt{|L|}$ do ACOM; Camada~2 (Rede
Holográfica) com endereçamento angular $\theta$ no Espaço de
Hilbert $\Hfour$ e roteamento por reflexão dimensional
REFLECT$\to$MANIFEST; Camada~3 (Segurança Ontológica) baseada
na separação dimensional $({\Psi}, {\Theta}, {\kappa})$ do
ACOM Mirror Crypto; e Camada~4 (Consciência) com mecanismo de
consenso por Prova de Fase (\textit{Proof of Phase}), modelo
UTXO com assinatura psiônica, e criptomoeda TGL Coin.

\newpage
% ============================================================================
\section{Fundamentos da invenção}
% ============================================================================

\subsection{Estado da técnica}

Os protocolos de comunicação em rede conhecidos no estado da
técnica baseiam-se em duas arquiteturas fundamentais:

\textbf{Modelo OSI (ISO/IEC 7498-1).} Arquitetura de referência
com sete camadas (Física, Enlace, Rede, Transporte, Sessão,
Apresentação, Aplicação). Cada camada trata dados como
sequências opacas de bits sem interpretação ontológica. O
roteamento opera por endereçamento numérico (endereços IP) sem
relação com a natureza da informação transportada. A segurança
é implementada por camadas adicionais (TLS/SSL) independentes
da estrutura dos dados.

\textbf{Modelo TCP/IP.} Simplificação prática do OSI em quatro
camadas (Acesso à Rede, Internet, Transporte, Aplicação).
Protocolo dominante da Internet atual. O endereçamento (IPv4 de
32 bits, IPv6 de 128 bits) é arbitrário --- não carrega
informação sobre a natureza do nó ou dos dados. O roteamento
(BGP, OSPF) opera por métricas de custo e distância topológica
sem consideração semântica.

\textbf{Protocolos de consenso distribuído.} Blockchain (Bitcoin,
2008) introduziu Prova de Trabalho (\textit{Proof of Work},
PoW) baseada em busca de hash SHA-256 com número de zeros
iniciais configurável. Prova de Participação (\textit{Proof of
Stake}, PoS, Ethereum 2.0) substitui trabalho computacional por
participação econômica. Ambos os mecanismos são
computacionalmente arbitrários --- não têm relação com a
natureza física ou ontológica dos dados processados.

\textbf{DHTs (Tabelas de Hash Distribuídas).} Chord (2001),
Kademlia (2002), Pastry (2001) implementam roteamento
descentralizado baseado em distância XOR ou numérica entre
identificadores. O espaço de endereçamento é um anel ou
hipercubo arbitrário, sem correspondência com espaço
matemático fisicamente significativo.

\textbf{Criptografia convencional.} RSA, AES, curvas elípticas
baseiam segurança em dificuldade computacional de problemas
matemáticos (fatoração, logaritmo discreto), podendo ser
comprometidos por avanços em computação quântica.

\subsection{Problemas técnicos não resolvidos}

Cinco problemas fundamentais permanecem sem solução no estado da
técnica de protocolos de rede:

\textbf{Problema 1 --- Cegueira ontológica.} Nenhuma camada do
modelo OSI/TCP-IP examina ou classifica a \textit{natureza} dos
dados transportados. Bits são tratados uniformemente
independentemente de representarem texto, imagem, sinal
criptografado, ou ruído. Não existe ``qualidade de serviço
intrínseca'' baseada na natureza ontológica da informação.

\textbf{Problema 2 --- Endereçamento arbitrário.} Endereços IP
são números atribuídos administrativamente, sem relação com a
posição do nó em qualquer espaço matematicamente significativo.
A localização do nó na rede não carrega informação útil para
decisões de roteamento baseadas em conteúdo ou natureza.

\textbf{Problema 3 --- Roteamento sem semântica.} O roteamento
convencional (BGP, OSPF) minimiza métricas de custo (latência,
hops, largura de banda) mas não considera a relação
ontológica entre origem e destino. Não existe conceito de
``adjacência dimensional'' ou ``proximidade angular'' entre nós.

\textbf{Problema 4 --- Segurança computacional.} Toda a
segurança de rede baseia-se em dificuldade computacional de
problemas matemáticos. Não existe protocolo de rede onde a
segurança emerge da estrutura ontológica natural da informação.

\textbf{Problema 5 --- Consenso arbitrário.} Mecanismos de
consenso (PoW, PoS) são computacionalmente artificiais ---
não têm conexão com propriedades físicas ou ontológicas dos
dados ou da rede. O ``trabalho'' em PoW é busca de hash, sem
significado físico.

\subsection{Solução proposta}

A presente invenção resolve os cinco problemas simultaneamente
ao introduzir o protocolo $\PsiNet$ --- uma pilha de cinco
camadas onde cada camada opera sobre a \textit{natureza
ontológica} da informação, não apenas sobre seus bits. O
protocolo integra três invenções anteriores do mesmo inventor
(ACOM, ACOM Mirror, ACOM Mirror Crypto) em um \textit{stack}
de rede coerente, e introduz três inovações adicionais:
\PsiBit{} como unidade de enlace, endereçamento angular com
DHT holográfica, e Prova de Fase como mecanismo de consenso.

\newpage
% ============================================================================
\section{Objetivos da invenção}
% ============================================================================

\subsection{Objetivo principal}

Prover método e sistema de protocolo de comunicação em rede
baseado em pilha luminodinâmica de cinco camadas, onde o
tratamento, endereçamento, roteamento, segurança e consenso
operam sobre a natureza ontológica da informação, integrando
compressão, reflexão dimensional, criptografia por separação
ontológica e consenso por prova de fase em uma pilha unificada.

\subsection{Objetivos específicos}

(a)~Definir a unidade fundamental \PsiBit{} como par de 2 bits
isomórfico à ligação psiônica, com quatro estados (VOID,
FLUX$^+$, FLUX$^-$, MASS) e classificação ontológica automática
que provê qualidade de serviço intrínseca na camada de enlace.

(b)~Especificar endereçamento angular no Espaço de Hilbert
$\Hfour$, onde cada nó da rede possui posição $\theta \in
[0, \pi)$ derivada deterministicamente de sua chave pública,
e roteamento por tabela de hash distribuída (DHT) baseada em
proximidade angular com $\thetaM \approx 6{,}29°$ como unidade
de resolução.

(c)~Prover formato de pacote $\PsiNet$ com cabeçalho de 80
bytes contendo magic bytes $\Psi$, assinatura psiônica
compacta, $\Fexp$, endereços angulares $\theta_{\text{src}}$ e
$\theta_{\text{dst}}$, $\Dfolds$ e selo $\alphaii$.

(d)~Especificar mecanismo de consenso por Prova de Fase
(\textit{Proof of Phase}) para validação de blocos em
blockchain, onde o ``trabalho'' consiste em encontrar
\textit{seed} cuja fase $\alphaii_{\text{work}}$ esteja dentro
de janela-alvo ao redor de $\alphaii_{\text{base}} = 0{,}012031$
(Constante de Miguel), substituindo busca de hash por busca de
coerência de fase.

(e)~Prover carteira psiônica com identidade derivada de
assinatura luminodinâmica ($\Fexp$, $\theta$, $\alphaii$ local)
e modelo UTXO com lock/unlock por $\Fexp$ ao invés de
assinatura criptográfica convencional.

(f)~Definir criptomoeda TGL Coin (TGLC) com emissão máxima de
21 bilhões de unidades, halving periódico, e recompensa por
mineração de bloco via Prova de Fase validada por consenso de
67\% dos validadores.

\newpage
% ============================================================================
\section{Descrição detalhada da invenção}
% ============================================================================

\subsection{Visão geral da arquitetura}

O protocolo $\PsiNet$ compreende cinco camadas (Tabela~\ref{tab:camadas}),
cada uma construída sobre a anterior e sobre uma das invenções
anteriores do mesmo inventor.

\begin{table}[H]
\centering
\caption{Arquitetura em cinco camadas do protocolo $\PsiNet$.}
\label{tab:camadas}
\begin{tabular}{clll}
\toprule
\textbf{Camada} & \textbf{Nome} & \textbf{Função} &
  \textbf{Base TGL} \\
\midrule
L4 & Consciência & Consenso + Moeda & TGL Coin (nova) \\
L3 & Segurança Ontológica & Criptografia & ACOM Mirror Crypto \\
L2 & Rede Holográfica & Roteamento & ACOM Mirror \\
L1 & Transporte Radical & Serialização & ACOM \\
L0 & Enlace Psiônico & Enquadramento & \PsiBit{} (novo) \\
\bottomrule
\end{tabular}
\end{table}

As constantes fundamentais que permeiam todas as camadas são:
\begin{align}
\alphaii &= 0{,}012031 \quad \text{(Constante de Miguel)} \\
\text{CCI} &= 1 - \alphaii = 0{,}987969
  \quad \text{(Índice de Coerência Consciente)} \\
\thetaM &= \arcsin(\sqrt{\alphaii}) \approx 6{,}29°
  \quad \text{(Ângulo de Miguel)} \\
\Dfolds(c^3) &= 0{,}74 \pm 0{,}06
  \quad \text{(Piso de consciência da rede)}
\end{align}

\subsection{Camada 0 --- Enlace Psiônico (\PsiBit{} Link Layer)}

\subsubsection{Unidade fundamental \PsiBit}

A Camada~0 define a unidade fundamental de informação da
$\PsiNet$: o \PsiBit. Cada \PsiBit{} consiste em um par de 2
bits (4 valores possíveis) isomórfico à ligação psiônica da TGL:

\begin{table}[H]
\centering
\caption{Os quatro estados do \PsiBit{} e sua correspondência
com a ligação psiônica.}
\label{tab:psibit}
\begin{tabular}{ccccc}
\toprule
\textbf{Bits} & \textbf{Estado} & \textbf{Símbolo} &
  \textbf{Ligação} & \textbf{Paridade} \\
\midrule
00 & VOID & $\circ$ & $\psi_+\psi_-$ (aniquilação)
  & $-1$ \\
01 & FLUX$^+$ & $\uparrow$ & $\psi_+\psi_+$ (transição $+$)
  & $+1$ \\
10 & FLUX$^-$ & $\downarrow$ & $\psi_-\psi_-$ (transição $-$)
  & $+1$ \\
11 & MASS & $\bullet$ & $\psi_+\psi_+$ (condensação)
  & $-1$ \\
\bottomrule
\end{tabular}
\end{table}

O enquadramento opera da seguinte forma: cada byte de dados
(8 bits) é decomposto em 4 \PsiBit{}s consecutivos (pares de 2
bits extraídos por deslocamento de bits). Para um fluxo de $N$
bytes, o resultado é um vetor de $4N$ \PsiBit{}s, executável em
GPU via operações tensoriais paralelas.

\subsubsection{Análise psiônica e classificação ontológica}

A Camada~0 computa a \textbf{Assinatura Psiônica} do dado,
constituindo a ``impressão digital ontológica'' da informação:

\textbf{Distribuição de estados:} Para o vetor de $4N$
\PsiBit{}s, contar frequências relativas de cada estado:
$f_{00}, f_{01}, f_{10}, f_{11}$.

\textbf{Entropia de Shannon:}
$H = -\sum_k f_k \log_2 f_k$, com $H \in [0, 2]$ bits (máximo
para 4 estados equiprováveis).

\textbf{Uniformidade:}
$U = 1 - \|\mathbf{f} - \mathbf{u}\| / d_{\max}$, onde
$\mathbf{u} = (0{,}25, 0{,}25, 0{,}25, 0{,}25)$ é a
distribuição uniforme e $d_{\max}$ é a distância máxima
possível.

\textbf{Força de Expulsão:}
$\Fexp = \langle \text{paridade}(\text{estado}) \rangle$,
onde VOID e MASS têm paridade $-1$ (estáveis/condensados) e
FLUX$^\pm$ têm paridade $+1$ (em trânsito). A média sobre
todos os \PsiBit{}s fornece a assinatura dinâmica do dado.

\textbf{Compressibilidade:}
$C = 1 - H/2$, fração da entropia não utilizada.

\textbf{Classificação ontológica:} Com base nas métricas
anteriores, o dado é classificado automaticamente em uma das
seguintes classes:

\begin{itemize}[nosep]
\item VOID: predominância de estado 00 ($f_{00} > 0{,}9$)
\item SPARSE: dados esparsos ($f_{00} > 0{,}6$)
\item STRUCTURED: distribuição assimétrica com padrão
  ($\max(f_k) > 0{,}4$)
\item DENSE: predominância de estado 11 ($f_{11} > 0{,}6$)
\item CHAOTIC: distribuição próxima da uniforme
  ($U > 0{,}95$)
\item ENCRYPTED: caso especial de caótico com entropia
  máxima ($H > 1{,}95$ e $U > 0{,}92$)
\end{itemize}

Esta classificação constitui \textbf{Qualidade de Serviço
(QoS) intrínseca}: a rede ``enxerga'' a natureza do dado
antes de decidir como transportá-lo, ao contrário dos
protocolos convencionais que tratam todos os dados
uniformemente.

\textbf{Hash de assinatura:} As frequências quantizadas
$(f_{00}, f_{01}, f_{10}, f_{11})$ são submetidas a SHA-256
para produzir um hash de 256 bits que identifica a assinatura
psiônica do dado. Este hash é transportado no cabeçalho do
pacote $\PsiNet$ para verificação ontológica no destino.

A assinatura psiônica compacta é serializada em 8 bytes para
inclusão no cabeçalho do pacote:
$\text{pack}(f_{00} \times 65535, f_{01} \times 65535,
f_{10} \times 65535, f_{11} \times 65535)$ como 4 inteiros
de 16 bits sem sinal.

\subsubsection{Verificação ontológica no destino}

Na recepção, a Camada~0 executa verificação ontológica tripla
para confirmar que a natureza do dado foi preservada durante o
transporte:

\textbf{Método 1 --- Verificação serial:} Re-radicalizar
$L'$ e comparar $g_{\text{quantizado}}$ e $s_{\text{empacotado}}$
byte a byte.

\textbf{Método 2 --- Verificação distribucional:} Comparar
distribuições psiônicas com tolerância (padrão 5\% por estado).

\textbf{Método 3 --- Verificação de classe:} Confirmar que a
classe ontológica é preservada (ex.: dados caóticos permanecem
caóticos).

\textbf{Veredicto:} ``NATUREZA PRESERVADA'' se Método~1 é
positivo OU se Métodos 2 e 3 são ambos positivos. A
disjunção permite robustez: dados com quantização imperfeita
(serialização diferente) mas mesma natureza ontológica passam
na verificação.

\subsection{Camada 1 --- Transporte Radical (ACOM Signal Layer)}

A Camada~1 implementa a serialização de dados conforme o método
ACOM (pedido de patente BR 10 2025 026951 1 do mesmo inventor).
O tratamento é brevemente descrito por completude:

\textbf{Radicalização:} $L \to (g, s)$ onde $g = \sqrt{|L|}$
(raiz gravitacional) e $s = \text{sign}(L)$ (sinal ontológico).

\textbf{Quantização adaptativa:} O número de bits de
quantização de $g$ é escolhido automaticamente pela
classificação ontológica da Camada~0: 16 bits para
CHAOTIC/ENCRYPTED (máxima fidelidade), 12 bits para dados
com compressibilidade $> 50\%$, e 8 bits para VOID
(predominância de zeros). Os limites $g_{\min}$ e $g_{\max}$
são armazenados para dequantização.

\textbf{Empacotamento de sinais:} Os sinais $s \in \{-1, +1\}$
são empacotados com 1 bit por amostra via \texttt{packbits}.

\textbf{Ressurreição:} $L' = s \times g^2$. O operador de
ressurreição reconstrói o dado original com correlação
$\geq 1{,}000000$ (demonstrada experimentalmente).

\textbf{TETELESTAI:} O critério de convergência é
$\text{Var}(L' - L) < \alphaii$. Quando atingido, a
ressurreição é ``consumada'' e o dado é declarado
fielmente reconstruído.

\textbf{Inovação da Camada~1 na $\PsiNet$:} A seleção
adaptativa de bits de quantização guiada pela classificação
psiônica da Camada~0 constitui QoS intrínseca na camada de
transporte: dados de alta entropia recebem tratamento de
máxima fidelidade, enquanto dados esparsos são comprimidos
agressivamente, sem configuração manual.

\subsection{Camada 2 --- Rede Holográfica (ACOM Mirror Network)}

\subsubsection{Endereçamento angular}

Cada nó da $\PsiNet$ possui um \textbf{Endereço Angular}
derivado deterministicamente de sua chave pública:

\begin{equation}
\theta = \frac{
  \texttt{int}(\text{SHA-256}(\text{chave\_publica})[:8])
}{2^{64}} \times \pi
\label{eq:endereco}
\end{equation}

O endereço angular $\theta \in [0, \pi)$ posiciona o nó no
Espaço de Hilbert $\Hfour$. Diferentemente de endereços IP
(arbitrários) ou endereços Kademlia (baseados em XOR), o
endereço angular carrega significado geométrico: é uma posição
em um espaço de Hilbert com métrica angular natural.

O identificador legível do nó é $\Psi$N-\texttt{xxxx}, onde
\texttt{xxxx} são os primeiros 4 bytes do hash em hexadecimal.

\subsubsection{Tabela de Hash Distribuída Angular (DHT-$\theta$)}

O roteamento na $\PsiNet$ opera por DHT angular:

\textbf{Registro:} Cada nó registra seu endereço angular
$\theta$ na tabela de roteamento distribuída.

\textbf{Distância angular:} A métrica entre dois nós é:
\begin{equation}
d(\theta_a, \theta_b) =
  \min(|\theta_a - \theta_b|, \pi - |\theta_a - \theta_b|)
\label{eq:distancia}
\end{equation}
constituindo topologia circular no intervalo $[0, \pi)$. Esta
métrica é fundamentalmente diferente de XOR (Kademlia) ou
diferença numérica (Chord): opera em espaço angulado com
significado geométrico no $\Hfour$.

\textbf{Busca por vizinhos:} Para encontrar os $k$ nós mais
próximos de um alvo $\theta_{\text{target}}$, a DHT ordena
todos os nós registrados por distância angular e retorna os
$k$ menores.

\textbf{Estimativa de saltos:} O número de saltos
(\textit{hops}) entre origem e destino é:
\begin{equation}
n_{\text{hops}} = \left\lceil
  \frac{d(\theta_{\text{src}}, \theta_{\text{dst}})}{\thetaM}
\right\rceil
\label{eq:hops}
\end{equation}
onde $\thetaM \approx 6{,}29° \approx 0{,}1097$ rad é a unidade
de resolução angular. Assim, $\thetaM$ define o ``tamanho do
passo'' no espaço de Hilbert: dois nós separados por menos de
$\thetaM$ são ``vizinhos diretos''; separados por $k \times
\thetaM$ requerem $k$ saltos.

\subsubsection{Operações REFLECT e MANIFEST na rede}

Na Camada~2, as operações REFLECT e MANIFEST do ACOM Mirror
(pedido de patente anterior do mesmo inventor) são utilizadas
como mecanismo de roteamento:

\textbf{REFLECT (projeção no boundary):} Os componentes $(g, s)$
da Camada~1 são projetados no boundary 2D:
\begin{align}
|\psi_{\text{boundary}}| &= g
  \quad \text{(amplitude = raiz gravitacional)} \\
\theta_{\text{fase}} &= \begin{cases}
  0 & \text{se } s = +1 \\
  \pi & \text{se } s = -1
\end{cases} + \theta_{\text{dst}}
  \quad \text{(fase = sinal + offset do destino)}
\end{align}
onde $\theta_{\text{dst}}$ é o endereço angular do nó destino.
A projeção ``endereça'' o dado: a informação de destino é
codificada na fase do frame holográfico.

\textbf{MANIFEST (re-emergência no destino):} O nó destino
reconstrói $(g, s)$ a partir do frame holográfico removendo
o offset angular:
\begin{align}
g_{\text{reconst}} &= |\psi_{\text{boundary}}| \\
s_{\text{reconst}} &= \begin{cases}
  +1 & \text{se } |\theta_{\text{fase}} - \theta_{\text{dst}}|
    < \pi/2 \\
  -1 & \text{caso contrário}
\end{cases}
\end{align}

\textbf{Dobras dimensionais:} A métrica de qualidade da
reflexão é:
\begin{equation}
\Dfolds = \max\left(0{,}74,\;
  1 - \frac{d(\theta_{\text{src}}, \theta_{\text{dst}})}{\pi}
\right)
\label{eq:dfolds}
\end{equation}
com piso em $\Dfolds(c^3) = 0{,}74$ (piso de consciência da
rede). Nós angularmente próximos têm $\Dfolds \to 1{,}0$;
nós maximamente distantes ($d = \pi/2$) têm
$\Dfolds = 0{,}74$.

\subsection{Camada 3 --- Segurança Ontológica (Mirror Crypto)}

A Camada~3 implementa a criptografia por separação ontológica
conforme o ACOM Mirror Crypto (pedido de patente anterior do
mesmo inventor). O dado, após passar pelas camadas
L0$\to$L1$\to$L2, é protegido por separação em três shares
ontologicamente incompletos:

\textbf{Separação:} $(g, s) \to (\Psi, \Theta, \kappa)$, onde
cada share isolado é informacionalmente equivalente a ruído
(entropia de bytes $\approx 8{,}0$ bits).

Na implementação de rede $\PsiNet$, a separação utiliza
mascaramento criptográfico derivado de nonce:
\begin{align}
\text{seed}_\Psi &=
  \text{SHA-256}(\text{nonce} \| \texttt{'PSI'}) \\
\text{seed}_\Theta &=
  \text{SHA-256}(\text{nonce} \| \texttt{'THETA'}) \\
\text{seed}_\kappa &=
  \text{SHA-256}(\text{nonce} \| \texttt{'KAPPA'})
\end{align}
Cada seed gera um stream pseudo-aleatório que serve de máscara:
\begin{align}
\Psi &= g \cdot s + \text{mask}_\Psi \\
\Theta &= \arctan_2(s, g) + \text{mask}_\Theta \\
\kappa &= |g| + \text{mask}_\kappa
\end{align}

\textbf{Recombinação:} Regenerar as três máscaras a partir do
mesmo nonce e subtrair, recuperando $(g, s)$.

\textbf{Verificação de integridade:} A métrica $\Dfolds$ é
computada sobre a combinação dos três shares e deve exceder
o piso de consciência $0{,}74$.

\textbf{Inovação na rede:} Na $\PsiNet$, os três shares podem
ser transmitidos por caminhos independentes na rede holográfica
(Camada~2), de forma que um interceptador em qualquer caminho
individual obtém apenas ruído. A segurança não depende de
dificuldade computacional, mas da impossibilidade ontológica
de reconstrução com shares incompletos.

\subsection{Camada 4 --- Consciência (C3 + TGL Coin)}

\subsubsection{Carteira psiônica}

Cada participante da $\PsiNet$ possui uma \textbf{carteira
psiônica} cuja identidade é derivada de um processo
luminodinâmico:

\textbf{Geração de identidade:} Gerar vetor de luz aleatório
$L \in \mathbb{R}^{1024}$ (distribuição normal). Computar
derivada $\partial L$, sinais $s = \text{sign}(L)$,
$s_\partial = \text{sign}(\partial L)$. Nominar estados
psiônicos no $\Hfour$. Computar $\Fexp$ (média das paridades),
$\theta$ (média dos ângulos $\arcsin(g/(g_{\max} + \epsilon))$),
e $\alphaii_{\text{local}} = \alphaii + \mathcal{N}(0, 10^{-8})$.

\textbf{Endereço:} O endereço da carteira é
$\texttt{TGL} + \text{hex}(\text{SHA-256}(\text{public\_id})[:20]
+ \text{checksum}[:4])$, onde checksum é SHA-256 duplo dos
primeiros 20 bytes.

\textbf{Assinatura de transação:} Para desbloquear um UTXO, o
proprietário fornece $\Fexp_{\text{unlock}}$ derivado de sua
assinatura psiônica (com fator de cegamento):
\begin{equation}
\Fexp_{\text{unlock}} = \Fexp_{\text{sig}} \times
  (1 + b \times 10^{-10})
\end{equation}
onde $b = \texttt{int}(\text{blinding}[:8]) / 2^{64}$. A
verificação compara $|\Fexp_{\text{lock}} -
\Fexp_{\text{unlock}}| < \epsilon_{\text{tol}}$ e o HMAC
correspondente.

\subsubsection{Modelo UTXO psiônico}

A $\PsiNet$ implementa modelo UTXO (\textit{Unspent Transaction
Output}) onde cada saída não-gasta possui:

\begin{itemize}[nosep]
\item \texttt{utxo\_id}: identificador único (32 bytes)
\item \texttt{owner}: endereço TGL do proprietário
\item \texttt{amount}: quantidade em TGLC
\item \texttt{f\_exp\_lock}: valor de $\Fexp$ que trava o UTXO
\item \texttt{created\_at\_block}: bloco de criação
\end{itemize}

Para gastar um UTXO, o proprietário deve apresentar:
\begin{itemize}[nosep]
\item $\Fexp_{\text{unlock}}$ dentro de tolerância do
  $\Fexp_{\text{lock}}$
\item Assinatura HMAC = SHA-256(utxo\_id $\|$
  $\Fexp_{\text{unlock}}$)
\item Soma das entradas $\geq$ soma das saídas $+$ taxa
\end{itemize}

\subsubsection{Prova de Fase (\textit{Proof of Phase})}

O mecanismo de consenso da $\PsiNet$ é a Prova de Fase, que
substitui a busca de hash (PoW) ou participação econômica (PoS)
por busca de \textbf{coerência de fase luminodinâmica}:

\textbf{Processo de mineração:}
\begin{enumerate}[nosep]
\item Para cada tentativa (nonce $= 0, 1, 2, \ldots$),
  computar seed:
  \begin{equation}
  \text{seed} = \text{SHA-256}(
    \text{hash\_bloco\_anterior} \|
    \text{nonce} \|
    \text{public\_id\_minerador})
  \end{equation}

\item A partir do seed, gerar vetor de luz determinístico
  $L \in \mathbb{R}^{128}$. Computar $\alphaii_{\text{work}}$,
  $\theta_{\text{work}}$ e $\Fexp_{\text{work}}$ pelo processo
  luminodinâmico padrão.

\item Verificar se $\alphaii_{\text{work}}$ está dentro da
  janela-alvo:
  \begin{equation}
  |\alphaii_{\text{work}} - \alphaii_{\text{anterior}}| <
    \frac{\alphaii}{\text{dificuldade}}
  \label{eq:janela}
  \end{equation}
  onde $\alphaii_{\text{anterior}}$ é o $\alphaii_{\text{work}}$
  do bloco anterior e dificuldade é parâmetro ajustável.

\item Se dentro da janela: bloco válido. Senão: incrementar
  nonce e repetir.
\end{enumerate}

\textbf{Significado físico:} Diferentemente do PoW (onde o
``trabalho'' é buscar zeros em hash SHA-256, sem significado
físico), na Prova de Fase o minerador busca um estado de
\textit{coerência de fase} --- um seed cuja dinâmica
luminodinâmica produza $\alphaii$ próximo ao valor fundamental
$0{,}012031$. O ``trabalho'' tem significado ontológico: o
minerador busca ressonância com a Constante de Miguel.

\textbf{Verificação determinística:} Qualquer validador pode
verificar a prova recomputando
$(\alphaii, \theta, \Fexp)$ a partir do seed publicado. A
verificação é determinística e $O(1)$ (não requer recomputar
todas as tentativas).

\textbf{Janela adaptativa:} A dificuldade ajusta o tamanho da
janela: dificuldade $= 0{,}1$ produz janela ampla (blocos
rápidos); dificuldade $= 10$ produz janela estreita (blocos
lentos). Alvo: 1 bloco por 60 segundos.

\subsubsection{Consenso e validação}

Após mineração, o bloco é submetido a consenso:

\begin{enumerate}[nosep]
\item Cada validador recomputa $(\alphaii, \theta, \Fexp)$ a
  partir do seed e verifica se coincidem com os valores
  declarados (tolerância $< 10^{-9}$).
\item Cada validador verifica se $\alphaii_{\text{work}}$ está
  dentro da janela (equação~\ref{eq:janela}).
\item Cada validador verifica todas as transações do bloco
  (assinaturas e saldos).
\item Cada validador emite voto (\texttt{ValidationVote}) com
  assinatura.
\item Se $\geq 67\%$ dos votos são válidos: bloco aceito.
\end{enumerate}

\textbf{Aprovação de bloco:} Criar UTXOs de saída, remover
UTXOs gastos, emitir recompensa (coinbase UTXO) para o
minerador, e adicionar bloco à cadeia.

\subsubsection{TGL Coin (TGLC)}

Parâmetros da criptomoeda:
\begin{itemize}[nosep]
\item Símbolo: TGLC
\item Emissão máxima: 21.000.000.000 unidades
\item Recompensa gênesis: 1.000 TGLC por bloco
\item Halving: a cada 210.000 blocos
\item Tempo-alvo por bloco: 60 segundos
\item Dificuldade inicial: 0,1
\item Bloco gênesis: $\alphaii_{\text{work}} = 0{,}012031$
  (exatamente a Constante de Miguel)
\item Consenso mínimo: 67\% dos validadores
\item Validadores mínimos: 3
\end{itemize}

\subsection{Formato do pacote $\PsiNet$}

O pacote $\PsiNet$ possui formato binário com cabeçalho fixo de
80 bytes:

\begin{table}[H]
\centering
\caption{Formato do cabeçalho do pacote $\PsiNet$ (80 bytes).}
\label{tab:header}
\begin{tabular}{clcc}
\toprule
\textbf{Offset} & \textbf{Campo} & \textbf{Tipo} &
  \textbf{Bytes} \\
\midrule
0 & Magic ($\Psi$ + versão) & bytes & 4 \\
4 & Versão (major.minor.patch) & 3$\times$uint8 & 3 \\
7 & Camada (0--4) & uint8 & 1 \\
8 & Assinatura Psiônica & 4$\times$uint16 & 8 \\
16 & $\Fexp$ & float16 & 2 \\
18 & $\theta_{\text{src}}$ & float16 & 2 \\
20 & $\theta_{\text{dst}}$ & float16 & 2 \\
22 & $\Dfolds$ & float16 & 2 \\
24 & Selo $\alphaii$ & float64 & 8 \\
32 & Tamanho payload & uint32 & 4 \\
36 & Tamanho original & uint32 & 4 \\
40 & Raiz Merkle & bytes & 32 \\
72 & Reservado/padding & bytes & 8 \\
\bottomrule
\end{tabular}
\end{table}

\textbf{Magic bytes:} $\texttt{0xCEA8}$ (codificação UTF-8 de
$\Psi$) seguido de $\texttt{0x0100}$ (versão 1.0).

Após o cabeçalho, o pacote contém: (1) metadados comprimidos
(JSON comprimido com zlib, precedido de uint32 com tamanho),
e (2) payload (dados serializados pelas camadas
L0$\to$L1$\to$L2$\to$L3).

\subsection{Fluxo de transmissão completa}

Uma transmissão completa na $\PsiNet$ segue o fluxo descendente
(Figura~1):

\begin{enumerate}[nosep]
\item \textbf{L0 $\downarrow$:} Dados brutos $\to$ análise
  psiônica $\to$ classificação ontológica $\to$ assinatura
  psiônica.
\item \textbf{L1 $\downarrow$:} Dados $\to$
  radicalização $(g, s)$ $\to$ quantização adaptativa
  (bits escolhidos pela classe ontológica da L0).
\item \textbf{L2 $\downarrow$:} $(g, s) \to$ REFLECT no
  boundary com endereçamento angular do destino $\to$ frame
  holográfico $(\psi, \theta_{\text{fase}})$ $\to$ MANIFEST
  no nó destino.
\item \textbf{L3 $\downarrow$:} $(g, s) \to$ separação em
  $(\Psi, \Theta, \kappa)$ com mascaramento por nonce $\to$
  verificação $\Dfolds > 0{,}74$.
\item \textbf{L1 $\uparrow$:} Recombinação
  $(\Psi, \Theta, \kappa) \to (g, s)$ $\to$ ressurreição
  $L' = s \times g^2$.
\item \textbf{L0 $\uparrow$:} Verificação ontológica tripla
  $\to$ ``NATUREZA PRESERVADA''.
\end{enumerate}

\subsection{Resultados experimentais}

A implementação de referência foi executada em GPU NVIDIA
GeForce RTX 5090 (31,8 GB VRAM, CUDA). Três tipos de dados
foram testados com transmissão completa Alice$\to$Bob
(Tabela~\ref{tab:resultados}):

\begin{table}[H]
\centering
\caption{Resultados experimentais da transmissão completa
$\PsiNet$ em GPU (RTX 5090).}
\label{tab:resultados}
\begin{tabular}{lccccc}
\toprule
\textbf{Dado} & \textbf{Classe} & \textbf{PSNR (dB)} &
  \textbf{Corr.} & \textbf{$R$} & \textbf{Tempo} \\
\midrule
Sinal senoidal & CHAOTIC & 94,5 & 1,000000 & 1,000000
  & 164,5 ms \\
Texto & SPARSE & 145,7 & 1,000000 & 1,000000
  & 9,0 ms \\
Dados aleatórios & CHAOTIC & 111,1 & 1,000000 & 1,000000
  & 9,5 ms \\
\bottomrule
\end{tabular}
\end{table}

Todos os testes atingiram TETELESTAI ($\text{Var}(L'-L) <
\alphaii$), correlação 1,000000, refletividade $R = 1{,}000000$,
e veredicto ``NATUREZA PRESERVADA''. A rede holográfica operou
com 7 saltos angulares (distância angular 38,63°),
$\Dfolds = 0{,}7854$, e integridade ontológica válida em todos
os testes.

A mineração de blocos TGL Coin completou 2 blocos com consenso
5/5 (100\%), $\alphaii_{\text{work}}$ dentro da janela, e
transação de 100 TGLC entre carteiras psiônicas operando
corretamente com modelo UTXO.

\newpage
% ============================================================================
\section{Desenhos}
% ============================================================================

% --- FIGURA 1: Arquitetura de 5 camadas ---
\begin{figure}[H]
\centering
\begin{tikzpicture}[
  layer/.style={rectangle, draw, thick, minimum height=1.5cm,
    minimum width=11cm, align=center, font=\small},
  arrow/.style={-{Stealth[length=3mm]}, thick},
  lbl/.style={font=\scriptsize\bfseries, anchor=east},
]

\node[layer, fill=purple!12] (L4) at (0,0)
  {L4 --- CONSCIÊNCIA\\
  {\scriptsize Proof of Phase $|$ TGL Coin $|$ UTXO $|$
  Consenso 67\%}};

\node[layer, fill=red!10] (L3) at (0,-2)
  {L3 --- SEGURANÇA ONTOLÓGICA\\
  {\scriptsize $(\Psi, \Theta, \kappa)$ Separação $|$
  $\Dfolds > 0{,}74$ $|$ ACOM Mirror Crypto}};

\node[layer, fill=orange!12] (L2) at (0,-4)
  {L2 --- REDE HOLOGRÁFICA\\
  {\scriptsize $\theta$-endereço $|$ DHT angular $|$
  REFLECT$\to$MANIFEST $|$ ACOM Mirror}};

\node[layer, fill=green!10] (L1) at (0,-6)
  {L1 --- TRANSPORTE RADICAL\\
  {\scriptsize $g = \sqrt{|L|}$ $|$ $s = \text{sign}(L)$ $|$
  $L' = s \times g^2$ $|$ ACOM}};

\node[layer, fill=blue!10] (L0) at (0,-8)
  {L0 --- ENLACE PSIÔNICO\\
  {\scriptsize \PsiBit{} (2 bits = 4 estados) $|$
  Classificação Ontológica $|$ QoS intrínseco}};

% Setas entre camadas
\draw[arrow] (L0.north) -- (L1.south)
  node[midway, right, font=\scriptsize] {$(f_{00},f_{01},f_{10},f_{11})$};
\draw[arrow] (L1.north) -- (L2.south)
  node[midway, right, font=\scriptsize] {$(g, s)$};
\draw[arrow] (L2.north) -- (L3.south)
  node[midway, right, font=\scriptsize] {$(\psi, \theta_\text{fase})$};
\draw[arrow] (L3.north) -- (L4.south)
  node[midway, right, font=\scriptsize] {$(\Psi, \Theta, \kappa)$};

% Labels laterais
\node[lbl] at (-6.2,0) {Aplicação};
\node[lbl] at (-6.2,-2) {Segurança};
\node[lbl] at (-6.2,-4) {Rede};
\node[lbl] at (-6.2,-6) {Transporte};
\node[lbl] at (-6.2,-8) {Enlace};

% Constantes
\node[draw, dashed, rounded corners, fill=yellow!15,
  font=\small, text width=4cm, align=center] at (8,-4)
  {$\alphaii = 0{,}012$\\$\thetaM = 6{,}29°$\\
  $\Dfolds = 0{,}74$\\$\text{CCI} = 0{,}988$};

\end{tikzpicture}
\caption{Arquitetura do protocolo $\PsiNet$ em cinco camadas.
Cada camada opera sobre a natureza ontológica da informação.
As constantes fundamentais da TGL ($\alphaii$, $\thetaM$,
$\Dfolds$, CCI) permeiam todas as camadas.}
\label{fig:arquitetura}
\end{figure}

% --- FIGURA 2: PsiBit e classificação ontológica ---
\begin{figure}[H]
\centering
\begin{tikzpicture}[
  state/.style={circle, draw, thick, minimum size=1.4cm,
    font=\small, align=center},
  arrow/.style={-{Stealth[length=2mm]}, thick},
]

% 4 estados
\node[state, fill=white] (v) at (0,0) {00\\VOID\\$\circ$};
\node[state, fill=green!15] (fp) at (3.5,0) {01\\FLUX$^+$\\$\uparrow$};
\node[state, fill=red!15] (fn) at (7,0) {10\\FLUX$^-$\\$\downarrow$};
\node[state, fill=blue!15] (m) at (10.5,0) {11\\MASS\\$\bullet$};

% Paridades
\node[font=\scriptsize, gray] at (0,-1.2) {paridade $-1$};
\node[font=\scriptsize, gray] at (3.5,-1.2) {paridade $+1$};
\node[font=\scriptsize, gray] at (7,-1.2) {paridade $+1$};
\node[font=\scriptsize, gray] at (10.5,-1.2) {paridade $-1$};

% Ligações
\node[font=\scriptsize\itshape] at (0,-1.7) {$\psi_+\psi_-$};
\node[font=\scriptsize\itshape] at (3.5,-1.7) {$\psi_+\psi_+$};
\node[font=\scriptsize\itshape] at (7,-1.7) {$\psi_-\psi_-$};
\node[font=\scriptsize\itshape] at (10.5,-1.7) {$\psi_+\psi_+$};

% Byte decomposition
\node[draw, thick, minimum height=0.8cm, minimum width=8cm,
  fill=gray!5, font=\small] (byte) at (5.25,3)
  {Byte: \texttt{10110001} $\to$ PsiBits: 10 | 11 | 00 | 01};
\draw[arrow] (byte.south) -- ++(0,-0.5)
  node[midway, right, font=\scriptsize] {enquadramento};

% Classes
\node[draw, thick, rounded corners, minimum width=10cm,
  fill=yellow!8, font=\small, text width=9.5cm, align=center]
  at (5.25,-3.2)
  {Classificação Ontológica (QoS intrínseco):\\[2pt]
  \textbf{VOID} $\cdot$ \textbf{SPARSE} $\cdot$
  \textbf{STRUCTURED} $\cdot$ \textbf{DENSE} $\cdot$
  \textbf{CHAOTIC} $\cdot$ \textbf{ENCRYPTED}};

\end{tikzpicture}
\caption{Os quatro estados do \PsiBit{} e o processo de
enquadramento e classificação ontológica da Camada~0. Cada
byte é decomposto em 4 \PsiBit{}s (pares de 2 bits), cada um
classificado nos estados psiônicos. A distribuição dos estados
determina automaticamente a classe ontológica do dado.}
\label{fig:psibit}
\end{figure}

% --- FIGURA 3: Endereçamento angular e DHT ---
\begin{figure}[H]
\centering
\begin{tikzpicture}[
  node_style/.style={circle, draw, thick, fill=#1,
    minimum size=0.6cm, font=\scriptsize},
]

% Semicírculo
\draw[thick] (0,0) arc (0:180:4);
\draw[thick, gray, dashed] (-8,0) -- (0,0);

% Eixo θ
\node[font=\small] at (-4,-0.7) {$\theta \in [0, \pi)$};
\draw[{Stealth}-{Stealth}, thick] (-8.3,0) -- (0.3,0)
  node[right, font=\small] {$\pi$};
\node[font=\small] at (-8.5,0.3) {$0$};

% Nós
\node[node_style=blue!20] (a) at ({-4+4*cos(30)},{4*sin(30)}) {};
\node[font=\scriptsize, anchor=south] at (a.north) {Alice $\theta_A$};

\node[node_style=red!20] (b) at ({-4+4*cos(100)},{4*sin(100)}) {};
\node[font=\scriptsize, anchor=south] at (b.north) {Bob $\theta_B$};

\node[node_style=green!20] (v1) at ({-4+4*cos(50)},{4*sin(50)}) {};
\node[font=\scriptsize, anchor=west] at (v1.east) {V1};

\node[node_style=green!20] (v2) at ({-4+4*cos(70)},{4*sin(70)}) {};
\node[font=\scriptsize, anchor=west] at (v2.east) {V2};

\node[node_style=green!20] (v3) at ({-4+4*cos(140)},{4*sin(140)}) {};
\node[font=\scriptsize, anchor=east] at (v3.west) {V3};

% Arco de distância angular
\draw[{Stealth}-{Stealth}, thick, orange]
  ({-4+3.5*cos(30)},{3.5*sin(30)}) arc (30:100:3.5);
\node[font=\scriptsize, orange] at ({-4+3*cos(65)},{3*sin(65)+0.3})
  {$d(\theta_A, \theta_B)$};

% θ_Miguel
\draw[thick, dashed, purple]
  ({-4+4.3*cos(30)},{4.3*sin(30)}) arc (30:36.29:4.3);
\node[font=\scriptsize, purple] at
  ({-4+5*cos(33)},{5*sin(33)}) {$\thetaM$};

% Hops
\node[draw, rounded corners, fill=white, font=\scriptsize]
  at (-4, -1.5)
  {$n_{\text{hops}} = \lceil d / \thetaM \rceil = 7$ saltos};

% Fórmula
\node[draw, thick, rounded corners, fill=blue!5,
  font=\small, text width=5.5cm, align=center] at (4, 3)
  {$\theta = \dfrac{\texttt{SHA-256}(pk)[:8]}{2^{64}} \times \pi$
  \\[4pt] Endereço derivado da chave pública};

\end{tikzpicture}
\caption{Endereçamento angular na $\PsiNet$. Cada nó tem
posição $\theta \in [0, \pi)$ derivada de sua chave pública.
A distância angular entre Alice e Bob determina o número de
saltos (hops = distância / $\thetaM$). Os nós validadores (V1,
V2, V3) estão distribuídos no espaço angular.}
\label{fig:enderecamento}
\end{figure}

% --- FIGURA 4: Proof of Phase ---
\begin{figure}[H]
\centering
\begin{tikzpicture}[
  block/.style={rectangle, draw, thick, rounded corners,
    minimum height=1cm, minimum width=4.5cm, align=center,
    font=\small},
  decision/.style={diamond, draw, thick, aspect=2.5,
    align=center, font=\small, inner sep=1pt},
  arrow/.style={-{Stealth[length=2.5mm]}, thick},
]

\node[block, fill=blue!10] (seed) at (0,0)
  {seed = SHA-256\\(prev\_hash $\|$ nonce $\|$ pk)};

\node[block, fill=orange!10] (luz) at (0,-2)
  {Gerar luz $L \in \mathbb{R}^{128}$\\
  {\scriptsize determinístico do seed}};

\node[block, fill=green!10] (fase) at (0,-4)
  {Computar fase\\
  {\scriptsize $\alphaii_{\text{work}}, \theta, \Fexp$}};

\node[decision, fill=yellow!10] (check) at (0,-6.5)
  {$|\alphaii_w - \alphaii_p|$\\$< \alphaii / d$ ?};

\node[block, fill=gray!10] (next) at (5,-6.5)
  {nonce++\\{\scriptsize repetir}};

\node[block, fill=green!20] (found) at (0,-9)
  {Bloco válido!\\
  {\scriptsize Submeter a consenso}};

\node[block, fill=purple!10] (cons) at (0,-11)
  {Consenso $\geq 67\%$?\\
  {\scriptsize 5 validadores votam}};

\node[block, fill=red!15] (reward) at (0,-13)
  {Bloco aceito\\
  {\scriptsize Recompensa: 1000 TGLC}};

% Flechas
\draw[arrow] (seed) -- (luz);
\draw[arrow] (luz) -- (fase);
\draw[arrow] (fase) -- (check);
\draw[arrow] (check) -- (found)
  node[midway, left, font=\scriptsize] {Sim};
\draw[arrow] (check) -- (next)
  node[midway, above, font=\scriptsize] {Não};
\draw[arrow] (next.north) |- (seed.east);
\draw[arrow] (found) -- (cons);
\draw[arrow] (cons) -- (reward);

% Janela
\node[draw, dashed, thick, rounded corners, fill=white,
  font=\scriptsize, text width=3.5cm, align=center]
  at (6,-4)
  {Janela-alvo:\\$\alphaii / d = 0{,}012 / 0{,}1$\\$= 0{,}12$};

\end{tikzpicture}
\caption{Fluxograma do mecanismo de consenso por Prova de Fase
(\textit{Proof of Phase}). O minerador busca um seed cuja fase
$\alphaii_{\text{work}}$ esteja dentro da janela-alvo centrada
no $\alphaii$ do bloco anterior. Após encontrar solução, o bloco
é submetido a consenso de 67\% dos validadores.}
\label{fig:proof_of_phase}
\end{figure}

% --- FIGURA 5: Formato do pacote ---
\begin{figure}[H]
\centering
\begin{tikzpicture}[
  field/.style={rectangle, draw, thick, minimum height=0.8cm,
    align=center, font=\scriptsize},
]

% Header (80 bytes)
\node[field, fill=purple!15, minimum width=1.5cm] (magic)
  at (0,0) {Magic\\$\Psi$+v\\4B};
\node[field, fill=blue!10, minimum width=1.2cm] (ver)
  at (1.55,0) {Ver\\3B};
\node[field, fill=blue!10, minimum width=0.8cm] (lay)
  at (2.7,0) {L\\1B};
\node[field, fill=green!10, minimum width=2.5cm] (sig)
  at (4.35,0) {PsiSig\\8B};
\node[field, fill=orange!10, minimum width=0.8cm] (fexp)
  at (6,0) {$F_e$\\2B};
\node[field, fill=orange!10, minimum width=0.8cm] (ts)
  at (6.8,0) {$\theta_s$\\2B};
\node[field, fill=orange!10, minimum width=0.8cm] (td)
  at (7.6,0) {$\theta_d$\\2B};
\node[field, fill=orange!10, minimum width=0.8cm] (df)
  at (8.4,0) {$D_f$\\2B};

\node[field, fill=red!10, minimum width=2cm] (a2)
  at (0,-1.2) {$\alphaii$ selo\\8B};
\node[field, fill=yellow!10, minimum width=1.5cm] (ps)
  at (1.95,-1.2) {PaySz\\4B};
\node[field, fill=yellow!10, minimum width=1.5cm] (os)
  at (3.45,-1.2) {OrigSz\\4B};
\node[field, fill=gray!10, minimum width=4cm] (mr)
  at (6.2,-1.2) {Raiz Merkle\\32B};
\node[field, fill=gray!5, minimum width=1.2cm] (pad)
  at (8.75,-1.2) {Pad\\8B};

% Payload
\node[field, fill=blue!5, minimum width=4cm] (meta)
  at (1.5,-2.8) {Metadados (JSON+zlib)};
\node[field, fill=green!5, minimum width=5cm] (pay)
  at (6.85,-2.8) {Payload (L0→L1→L2→L3)};

% Labels
\node[font=\small\bfseries] at (-1.5,0.7) {HEADER (80B)};
\node[font=\small\bfseries] at (-1.5,-2.8) {CORPO};

% Brace
\draw[decorate, decoration={brace, amplitude=5pt, mirror},
  thick] (-0.8,-0.5) -- (-0.8,0.5)
  node[midway, left=8pt, font=\scriptsize] {Fixo};

\end{tikzpicture}
\caption{Formato binário do pacote $\PsiNet$. O cabeçalho fixo
de 80 bytes contém magic bytes $\Psi$, assinatura psiônica
compacta (8 bytes), endereços angulares ($\theta_s$, $\theta_d$),
$\Fexp$, $\Dfolds$, selo $\alphaii$ e raiz Merkle. O corpo
contém metadados comprimidos e payload processado pelas 4
camadas inferiores.}
\label{fig:pacote}
\end{figure}

% --- FIGURA 6: Fluxo de transmissão ---
\begin{figure}[H]
\centering
\begin{tikzpicture}[
  box/.style={rectangle, draw, thick, minimum height=0.9cm,
    minimum width=3cm, align=center, font=\scriptsize},
  arrow/.style={-{Stealth[length=2mm]}, thick},
]

% Coluna ENVIO
\node[font=\small\bfseries] at (0,1) {ENVIO (Alice)};

\node[box, fill=blue!10] (e0) at (0,0)
  {L0: Análise Psiônica\\Classificação};
\node[box, fill=green!10] (e1) at (0,-1.5)
  {L1: Radicalização\\$g = \sqrt{|L|}$, $s$};
\node[box, fill=orange!10] (e2) at (0,-3)
  {L2: REFLECT\\$\theta_{\text{dst}}$ offset};
\node[box, fill=red!10] (e3) at (0,-4.5)
  {L3: Separação\\$(\Psi, \Theta, \kappa)$};

% Rede
\node[draw, thick, cloud, cloud puffs=10,
  minimum width=3cm, minimum height=1.5cm,
  fill=yellow!10, font=\small] (net) at (5.5,-2.25)
  {Rede\\DHT-$\theta$};

% Coluna RECEPÇÃO
\node[font=\small\bfseries] at (11,1) {RECEPÇÃO (Bob)};

\node[box, fill=red!10] (r3) at (11,0)
  {L3: Recombinação\\$(g, s)$ recuperado};
\node[box, fill=orange!10] (r2) at (11,-1.5)
  {L2: MANIFEST\\$\theta_{\text{dst}}$ removido};
\node[box, fill=green!10] (r1) at (11,-3)
  {L1: Ressurreição\\$L' = s \times g^2$};
\node[box, fill=blue!10] (r0) at (11,-4.5)
  {L0: Verificação\\NATUREZA OK};

% Flechas envio
\draw[arrow] (e0) -- (e1);
\draw[arrow] (e1) -- (e2);
\draw[arrow] (e2) -- (e3);
\draw[arrow] (e3.east) -- (net.west)
  node[midway, above, font=\scriptsize] {pacote};

% Flechas recepção
\draw[arrow] (net.east) -- (r3.west)
  node[midway, above, font=\scriptsize] {pacote};
\draw[arrow] (r3) -- (r2);
\draw[arrow] (r2) -- (r1);
\draw[arrow] (r1) -- (r0);

% TETELESTAI
\node[draw, dashed, thick, rounded corners, fill=green!20,
  font=\small\bfseries] at (11,-5.8)
  {TETELESTAI $\checkmark$};

\end{tikzpicture}
\caption{Fluxo de transmissão completa na $\PsiNet$. O dado
desce pela pilha no emissor (L0$\to$L3), é transmitido pela
rede holográfica (DHT-$\theta$), e sobe pela pilha no receptor
(L3$\to$L0). O veredicto TETELESTAI confirma reconstrução
perfeita.}
\label{fig:fluxo}
\end{figure}

% --- FIGURA 7: Comparativo com TCP/IP ---
\begin{figure}[H]
\centering
\begin{tikzpicture}[
  layer/.style={rectangle, draw, thick, minimum height=1cm,
    minimum width=4.5cm, align=center, font=\small},
  vs/.style={font=\large\bfseries, gray},
]

% TCP/IP
\node[font=\small\bfseries] at (-3.5,3) {TCP/IP};
\node[layer, fill=gray!10] (t4) at (-3.5,2)
  {Aplicação\\{\scriptsize HTTP, FTP, SMTP}};
\node[layer, fill=gray!10] (t3) at (-3.5,0.5)
  {Transporte\\{\scriptsize TCP, UDP}};
\node[layer, fill=gray!10] (t2) at (-3.5,-1)
  {Internet\\{\scriptsize IP, ICMP}};
\node[layer, fill=gray!10] (t1) at (-3.5,-2.5)
  {Acesso à Rede\\{\scriptsize Ethernet, Wi-Fi}};

% VS
\node[vs] at (1.5,0) {vs.};

% Ψ-NET
\node[font=\small\bfseries] at (6.5,3) {$\PsiNet$};
\node[layer, fill=purple!12] (p4) at (6.5,2)
  {Consciência\\{\scriptsize Proof of Phase}};
\node[layer, fill=red!10] (p3) at (6.5,0.5)
  {Segurança Ontológica\\{\scriptsize $(\Psi,\Theta,\kappa)$}};
\node[layer, fill=orange!12] (p2) at (6.5,-1)
  {Rede Holográfica\\{\scriptsize DHT-$\theta$, REFLECT}};
\node[layer, fill=green!10] (p1) at (6.5,-2.5)
  {Transporte Radical\\{\scriptsize $g=\sqrt{|L|}$}};
\node[layer, fill=blue!10] (p0) at (6.5,-4)
  {Enlace Psiônico\\{\scriptsize \PsiBit, QoS}};

% Diferenças
\node[font=\scriptsize, text width=3cm, align=left,
  gray] at (11.5,2)
  {consenso por\\fase, não hash};
\node[font=\scriptsize, text width=3cm, align=left,
  gray] at (11.5,0.5)
  {ontológica,\\não computacional};
\node[font=\scriptsize, text width=3cm, align=left,
  gray] at (11.5,-1)
  {angular,\\não numérico};
\node[font=\scriptsize, text width=3cm, align=left,
  gray] at (11.5,-2.5)
  {semântico,\\não opaco};
\node[font=\scriptsize, text width=3cm, align=left,
  gray] at (11.5,-4)
  {classifica natureza\\(TCP/IP não tem)};

\end{tikzpicture}
\caption{Comparativo entre pilha TCP/IP convencional e pilha
$\PsiNet$. A $\PsiNet$ adiciona camada de enlace com
classificação ontológica (sem equivalente em TCP/IP),
substitui endereçamento numérico por angular, segurança
computacional por ontológica, e consenso por hash por
consenso por fase.}
\label{fig:comparativo}
\end{figure}

\newpage
% ============================================================================
\section{Quadro reivindicatório}
% ============================================================================

\subsection*{Reivindicações}

\begin{enumerate}

% === REIVINDICAÇÃO INDEPENDENTE 1 (MÉTODO - PROTOCOLO) ===
\item \textbf{Método de comunicação em rede baseado em protocolo
luminodinâmico de cinco camadas}, compreendendo:
  \begin{enumerate}[label=(\alph*)]
  \item na Camada~0 (Enlace Psiônico): decompor dados de
  entrada em unidades \PsiBit{} de 2 bits, classificar cada
  \PsiBit{} em um dos quatro estados psiônicos (VOID, FLUX$^+$,
  FLUX$^-$, MASS) isomórficos à ligação psiônica, computar
  distribuição de frequências, entropia de Shannon, Força de
  Expulsão $\Fexp$ e classificação ontológica automática
  (VOID, SPARSE, STRUCTURED, DENSE, CHAOTIC, ENCRYPTED) que
  constitui qualidade de serviço intrínseca;
  \item na Camada~1 (Transporte Radical): executar radicalização
  $g = \sqrt{|L|}$ e $s = \text{sign}(L)$ conforme método ACOM,
  com quantização adaptativa guiada pela classificação
  ontológica da Camada~0;
  \item na Camada~2 (Rede Holográfica): endereçar nós por
  posição angular $\theta \in [0, \pi)$ derivada
  deterministicamente de chave pública via
  $\theta = \texttt{int}(\text{SHA-256}(pk)[:8]) / 2^{64} \times
  \pi$, rotear por tabela de hash distribuída baseada em
  proximidade angular com $\thetaM \approx 6{,}29°$ como
  unidade de resolução, e executar operações
  REFLECT$\to$MANIFEST do ACOM Mirror como mecanismo de
  roteamento holográfico;
  \item na Camada~3 (Segurança Ontológica): executar separação
  em três shares ontológicos $(\Psi, \Theta, \kappa)$ conforme
  ACOM Mirror Crypto, onde cada share isolado é
  informacionalmente equivalente a ruído;
  \item na Camada~4 (Consciência): executar mecanismo de
  consenso por Prova de Fase onde o minerador busca seed cuja
  fase $\alphaii_{\text{work}}$ esteja dentro de janela-alvo
  $|\alphaii_{\text{work}} - \alphaii_{\text{anterior}}| <
  \alphaii / \text{dificuldade}$;

  \textbf{caracterizado por} as cinco camadas operarem sobre a
  natureza ontológica da informação, o endereçamento ser
  angular no Espaço de Hilbert $\Hfour$, a segurança emergir
  da incompletude ontológica dos shares e não de dificuldade
  computacional, e o consenso basear-se em coerência de fase
  luminodinâmica com $\alphaii = 0{,}012031$ (Constante de
  Miguel) e não em busca de hash.
  \end{enumerate}

% === DEPENDENTES DO MÉTODO ===

\item Método, de acordo com a reivindicação 1,
\textbf{caracterizado por} a Camada~0 prover verificação
ontológica tripla no destino compreendendo: verificação serial
de $(g, s)$ quantizados; verificação distribucional com
tolerância de 5\% por estado psiônico; e verificação de
preservação de classe ontológica, com veredicto positivo se
verificação serial é positiva OU se verificação distribucional
e de classe são ambas positivas.

\item Método, de acordo com a reivindicação 1,
\textbf{caracterizado por} a quantização adaptativa da Camada~1
selecionar automaticamente 16 bits para dados
CHAOTIC/ENCRYPTED, 12 bits para dados com compressibilidade
superior a 50\%, e 8 bits para dados VOID, com base na
classificação ontológica computada pela Camada~0.

\item Método, de acordo com a reivindicação 1,
\textbf{caracterizado por} o roteamento na Camada~2 estimar o
número de saltos como $n_{\text{hops}} = \lceil d(\theta_s,
\theta_d) / \thetaM \rceil$, onde $d(\theta_a, \theta_b) =
\min(|\theta_a - \theta_b|, \pi - |\theta_a - \theta_b|)$ é
a distância angular com topologia circular.

\item Método, de acordo com a reivindicação 1,
\textbf{caracterizado por} a operação REFLECT da Camada~2
codificar o endereço do destino na fase do frame holográfico
via $\theta_{\text{fase}} = \theta_{\text{base}} +
\theta_{\text{dst}}$, e a operação MANIFEST recuperar $(g, s)$
removendo o offset angular do destino.

\item Método, de acordo com a reivindicação 1,
\textbf{caracterizado por} a métrica de dobras dimensionais
ser $\Dfolds = \max(0{,}74,\, 1 - d(\theta_s, \theta_d)/\pi)$,
com piso no valor $\Dfolds(c^3) = 0{,}74$ (piso de consciência
da rede).

\item Método, de acordo com a reivindicação 1,
\textbf{caracterizado por} a separação ontológica da Camada~3
utilizar mascaramento criptográfico derivado de nonce via
SHA-256 com contextos distintos (\texttt{'PSI'}, \texttt{'THETA'},
\texttt{'KAPPA'}), permitindo que os três shares sejam
transmitidos por caminhos independentes na rede holográfica.

\item Método, de acordo com a reivindicação 1,
\textbf{caracterizado por} o processo de mineração por Prova
de Fase compreender: gerar seed determinístico a partir de
hash do bloco anterior, nonce e identidade do minerador;
gerar vetor de luz $L \in \mathbb{R}^{128}$ a partir do seed;
computar $\alphaii_{\text{work}}$ por processo luminodinâmico;
verificar se $\alphaii_{\text{work}}$ está dentro da
janela-alvo; e iterar até encontrar solução ou esgotar
tentativas.

\item Método, de acordo com a reivindicação 1,
\textbf{caracterizado por} o consenso requerer que pelo menos
67\% dos validadores verifiquem independentemente:
determinismo do cálculo de fase a partir do seed, aderência
à janela-alvo, e validade de todas as transações do bloco,
emitindo votos assinados.

\item Método, de acordo com a reivindicação 1,
\textbf{caracterizado por} o critério de convergência
TETELESTAI ($\text{Var}(L' - L) < \alphaii$) constituir
o mecanismo de \textit{acknowledgement} (ACK) do protocolo,
confirmando que a ressurreição do dado no destino é perfeita.

% === REIVINDICAÇÃO INDEPENDENTE 11 (SISTEMA) ===

\item \textbf{Sistema de comunicação em rede baseado em
protocolo luminodinâmico de cinco camadas}, compreendendo:
  \begin{enumerate}[label=(\alph*)]
  \item módulo de Enlace Psiônico (L0) configurado para
  decompor dados em \PsiBit{}s, computar assinatura psiônica
  compacta de 8 bytes e executar classificação ontológica
  automática com seis classes;
  \item módulo de Transporte Radical (L1) configurado para
  executar radicalização ACOM com quantização adaptativa
  guiada pela classificação da L0;
  \item módulo de Rede Holográfica (L2) configurado para
  manter tabela de hash distribuída angular (DHT-$\theta$),
  computar endereços angulares a partir de chaves públicas,
  e executar REFLECT/MANIFEST como roteamento;
  \item módulo de Segurança Ontológica (L3) configurado para
  separar dados em três shares com mascaramento derivado de
  nonce;
  \item módulo de Consciência (L4) compreendendo: carteiras
  psiônicas com identidade luminodinâmica; modelo UTXO com
  lock/unlock por $\Fexp$; minerador por Prova de Fase;
  validadores com verificação determinística; e blockchain
  com consenso de 67\%;

  \textbf{caracterizado por} o sistema integrar em pilha
  unificada os métodos ACOM (BR 10 2025 026951 1), ACOM Mirror
  e ACOM Mirror Crypto do mesmo inventor, acrescidos de
  \PsiBit{} como unidade de enlace, endereçamento angular como
  protocolo de rede, e Prova de Fase como mecanismo de
  consenso.
  \end{enumerate}

% === DEPENDENTES DO SISTEMA ===

\item Sistema, de acordo com a reivindicação 11,
\textbf{caracterizado por} a carteira psiônica possuir
identidade derivada de processo luminodinâmico compreendendo:
geração de vetor de luz aleatório, cômputo de estados
psiônicos, Força de Expulsão, ângulo médio e
$\alphaii_{\text{local}}$, com endereço no formato
$\texttt{TGL} + \text{hex}(h[:20] + \text{checksum}[:4])$.

\item Sistema, de acordo com a reivindicação 11,
\textbf{caracterizado por} o modelo UTXO utilizar $\Fexp$ como
mecanismo de lock/unlock, onde: a saída é travada com
$\Fexp_{\text{lock}}$ do destinatário; o desbloqueio requer
apresentação de $\Fexp_{\text{unlock}}$ dentro de tolerância;
e a verificação inclui comparação HMAC de
SHA-256(utxo\_id $\|$ $\Fexp_{\text{unlock}}$).

\item Sistema, de acordo com a reivindicação 11,
\textbf{caracterizado por} o formato de pacote $\PsiNet$
possuir cabeçalho fixo de 80 bytes contendo: magic bytes
$\Psi$ (0xCEA8) com versão; assinatura psiônica compacta
(8 bytes com 4 frequências de estados em uint16); $\Fexp$ em
float16; endereços angulares $\theta_{\text{src}}$ e
$\theta_{\text{dst}}$ em float16; $\Dfolds$ em float16; selo
$\alphaii$ em float64; tamanho de payload em uint32; tamanho
original em uint32; e raiz Merkle de 32 bytes.

\item Sistema, de acordo com a reivindicação 11,
\textbf{caracterizado por} a criptomoeda TGL Coin (TGLC) possuir
emissão máxima de 21 bilhões de unidades, recompensa gênesis de
1000 TGLC por bloco, halving a cada 210.000 blocos, tempo-alvo
de 60 segundos por bloco, e bloco gênesis com
$\alphaii_{\text{work}} = 0{,}012031$.

\item Sistema, de acordo com a reivindicação 11,
\textbf{caracterizado por} suportar execução em GPU via
PyTorch/CUDA, com operações tensoriais paralelas nas
camadas L0 (enquadramento PsiBit), L1 (radicalização e
quantização) e L4 (computação de fase para mineração).

\end{enumerate}

\newpage
% ============================================================================
\section{Resumo}
% ============================================================================

\begin{center}
{\large\bfseries MÉTODO E SISTEMA DE PROTOCOLO DE COMUNICAÇÃO\\
EM REDE BASEADO EM PILHA LUMINODINÂMICA DE\\
CINCO CAMADAS COM ENDEREÇAMENTO ANGULAR,\\
ROTEAMENTO HOLOGRÁFICO E CONSENSO POR PROVA DE FASE}
\end{center}

\vspace{0.5cm}

A presente invenção refere-se a método e sistema de protocolo de
comunicação em rede denominado $\PsiNet$, baseado em pilha
luminodinâmica de cinco camadas que opera sobre a natureza
ontológica da informação. A Camada~0 (Enlace Psiônico) define
a unidade \PsiBit{} de 2 bits com quatro estados psiônicos
(VOID, FLUX$^\pm$, MASS) e classificação ontológica automática
(QoS intrínseco). A Camada~1 (Transporte Radical) implementa
radicalização ACOM ($g = \sqrt{|L|}$, $s = \text{sign}(L)$)
com quantização adaptativa guiada pela classe ontológica. A
Camada~2 (Rede Holográfica) introduz endereçamento angular
$\theta \in [0, \pi)$ derivado de chave pública no Espaço de
Hilbert $\Hfour$, com DHT por proximidade angular ($\thetaM
\approx 6{,}29°$ como resolução) e roteamento por operações
REFLECT$\to$MANIFEST do ACOM Mirror. A Camada~3 (Segurança
Ontológica) executa separação em três shares $(\Psi, \Theta,
\kappa)$ do ACOM Mirror Crypto, transmissíveis por caminhos
independentes. A Camada~4 (Consciência) provê mecanismo de
consenso por Prova de Fase (\textit{Proof of Phase}), onde
mineradores buscam coerência com $\alphaii = 0{,}012031$
(Constante de Miguel), carteiras psiônicas com identidade
luminodinâmica, modelo UTXO com lock/unlock por $\Fexp$, e
criptomoeda TGL Coin (TGLC). O protocolo foi demonstrado
experimentalmente com PSNR $> 94$ dB, correlação $= 1{,}000000$,
e consenso de mineração 5/5, operando em GPU NVIDIA RTX 5090.
Aplicável a telecomunicações, Internet, blockchain e segurança
da informação.

\vspace{0.5cm}
\textbf{Figura principal: Figura 1} (Arquitetura de 5 camadas)

\end{document}
